\section{Required Additional Work Before the Operator-Theoretic Claim}
\label{sec:requiredwork}

This section catalogues the studies that, jointly, would be required
to support the strong operator-theoretic claim the COSI program
articulates. None of these studies has been run in the present paper.
We list them explicitly, with the gap each addresses, so that the
present pilot result is read in proper context and so that subsequent
work has a clear program to execute.

\subsection{Study A. Invariance-Leakage Measurement}
\label{sec:study-a}

The paper defines a $k$-dimensional subspace $V \subseteq \R^d$ as
approximately $\epsilon$-invariant under a (possibly nonlinear) map
$f$ on a sample $X$ if
\begin{equation}
\frac{\fnorm{(I - P_V)\, f(P_V X)}}{\fnorm{f(P_V X)}} \;\leq\; \epsilon.
\label{eq:leakage}
\end{equation}
This is the central mathematical object the operator-theoretic framing
depends on. The pilot does not measure it. PCA top-$k$ components
identify high-variance directions; high-variance is not invariance.

\paragraph{Required.} For each model and each candidate subspace
identification method (PCA, sparse autoencoder features
\cite{bricken2023monosemanticity, cunningham2024sparse,
templeton2024scaling, gao2024scaling, lieberum2024gemma}, learned
linear probes \cite{hewitt2019structural, belinkov2022probing}),
project activations onto $V$, pass them through the next layer's
forward map, project the output onto $V$ and onto its orthogonal
complement, and compute the leakage statistic
\eqref{eq:leakage}. Report leakage as a function of layer, model,
domain, subspace dimension, and subspace-identification method. A
candidate subspace is operator-relevant only if its leakage is small
relative to its capacity to compress activations.

\paragraph{Gap addressed.} Whether the aligned subspaces are
operator-invariant in the sense the framing requires, or merely
high-variance principal components.

\subsection{Study B. Lexical, Template, and Syntactic Baselines}

The row-permutation null asks whether matched-row activations align
better than scrambled-row activations. This is a low bar. Two
competent language models exposed to the same English prompts will
preserve enough lexical and syntactic structure that matched-row
alignment exceeds scrambled-row alignment regardless of any
operator-theoretic substrate.

\paragraph{Required.} Run the same Procrustes alignment procedure on:

\begin{itemize}
    \item TF-IDF or bag-of-words representations of the same prompts
        in place of the model activations, as a lexical baseline.
    \item Sentence-transformer embeddings of the prompts as a
        learned-but-non-LLM baseline.
    \item Static word embeddings averaged over each prompt
        \cite{mikolov2013distributed} as a pre-transformer baseline.
    \item Prompts in which all content words have been replaced by
        placeholder tokens (template-only baseline).
    \item Prompts whose token sequence has been syntax-scrambled
        (preserving lexical content, destroying compositional order).
    \item Paraphrase-preserving variants where two semantically
        equivalent prompts replace each original prompt.
\end{itemize}

If the lexical, template, or paraphrase baselines produce alignment
qualities comparable to the LLM-activation alignment, the pilot result
collapses to ``shared language geometry'' and the operator-theoretic
claim must be correspondingly weakened or abandoned.

\paragraph{Gap addressed.} Whether the alignment isolates compositional
structure from lexical, template, syntactic, or tokenizer-induced
structure.

\subsection{Study C. Train/Test Procrustes and Cross-Domain Transfer}

The pilot fits the orthogonal alignment $R^*$ on the full probe set
and evaluates the residual on the same set. This conflates fit with
evaluation and provides no out-of-sample protection.

\paragraph{Required.} Fit $R^*$ on a randomly held-out training
fraction (e.g., $70\%$) of the probe set, evaluate the residual on
the held-out test fraction. Repeat across many random splits and
report the mean and confidence interval. More demanding: fit on one
domain (e.g., math) and evaluate on a different domain (e.g., logic),
testing whether the alignment is domain-general or domain-specific.

\paragraph{Gap addressed.} Whether the alignment is robust to held-out
prompts and whether it generalizes across domains rather than overfit
to the joint sample.

\subsection{Study D. Training-Distribution-Disjoint Controls}

Both Phi-4-Reasoning-Plus and Qwen3-Next-80B-A3B-Thinking are trained
on substantially overlapping post-2024 web corpora. The aligned
subspace could plausibly be a residue of shared training distribution
rather than a substrate-independent computational structure.

\paragraph{Required.} Replace one of the pilot pair's models with a
training-distribution-distant model, holding the other fixed.
Candidate disjoint-corpus controls:

\begin{itemize}
    \item A model trained primarily on a non-web scientific corpus,
        e.g., a model in the Galactica \cite{taylor2022galactica}
        lineage or a domain-specific scientific-literature model.
    \item A multilingual model whose training distribution is
        dominated by a non-English language.
    \item A small-scale model trained from scratch on a deliberately
        curated corpus with documented disjointness from the pilot
        pair's training data.
\end{itemize}

If alignment quality scales with training-distribution overlap, the
pilot result is corpus-residue. If alignment quality is preserved
across substantial training-distribution disjointness, the
substrate-independence reading is meaningfully supported.

\paragraph{Gap addressed.} Whether the alignment is a substrate-independent
computational structure or a shared training-corpus residue.

\subsection{Study E. Cross-Model Causal Intervention}

Procrustes alignment is correspondence evidence. The
operator-theoretic claim that the aligned subspace is causally
responsible for compositional reasoning behavior, rather than a
correlated statistical shadow, requires intervention.

\paragraph{Required.} Identify a direction $v_A \in V_A$ associated
with a specific compositional-reasoning behavior in $M_A$ via
established mechanistic-interpretability methods
\cite{wang2023interpretability, conmy2023towards, marks2024sparse,
meng2022rome}. Use the recovered Procrustes transformation $R^*$ to
predict the corresponding direction $v_B = R^{*\top} v_A \in V_B$ in
$M_B$. Predict the behavioral effect of ablating or steering along
$v_B$ in $M_B$. Run the intervention. Compare predicted to observed
effect across many such directions, reporting the prediction quality
relative to a control of randomly chosen directions in $V_B$.

\paragraph{Gap addressed.} Whether the alignment is causally
explanatory of compositional behavior in both models or merely
descriptive of activation geometry.

\subsection{Study F. Methodology Robustness Sweeps}

Several methodological choices in the pilot were not swept:

\begin{itemize}
    \item Pooling mode: last-token only; mean and max pooling and
        content-token / answer-position / relation-token pooling not
        reported.
    \item Layer fractions: only $0.25$, $0.5$, $0.75$. Layer-by-layer
        sweeps would resolve where the alignment is concentrated.
    \item PCA variance threshold: only $0.90$; sensitivity at $0.80$,
        $0.85$, $0.95$, $0.99$ unmeasured.
    \item Whitening: not tested. Procrustes residuals can be sensitive
        to differential variance structure across the top components
        of the two matrices.
    \item CKA \cite{kornblith2019similarity} cross-validation:
        promised in the methods but not yet computed and reported. We
        flag this as a methodological commitment that the pilot does
        not satisfy.
    \item SVCCA \cite{raghu2017svcca}, PWCCA \cite{morcos2018insights},
        and generalized shape metrics \cite{williams2021generalized}
        as further triangulation: not run.
    \item Effect-size reporting beyond standardized $z$-scores:
        confidence intervals, cluster bootstraps grouped by prompt
        template and vocabulary family, finite-permutation $p$-values
        with Monte Carlo standard errors.
\end{itemize}

\paragraph{Gap addressed.} Whether the pilot result is robust to
reasonable variation in the methodological choices that were fixed by
declaration rather than examined empirically.

\subsection{Study G. Externally Verifiable Pre-Registration}

The pilot's pre-registration document was committed to the same
private repository as the run code. This is not externally
auditable. Genuine pre-registration requires immutable timestamping
on a public registry.

\paragraph{Required.} For all subsequent COSI runs, deposit the
analysis plan, the exact prompt set with content hash, the model
checkpoint revisions, the random seeds, and the environment lockfile
on the Open Science Framework, Zenodo, or equivalent, with the
deposit timestamped before any data is observed. Cite the deposit
DOI in subsequent papers.

\paragraph{Gap addressed.} Whether the pre-registration commitments
the pilot states are externally verifiable rather than internal to a
private working repository.

\subsection{Summary: What Joint Completion of Studies A--G Would Establish}

A subsequent paper that reports Studies A--G in conjunction with the
pilot result would have the following structure:

\begin{itemize}
    \item The aligned subspaces are demonstrably approximately invariant
        under the forward operator (Study A), justifying the
        operator-theoretic interpretation of the alignment.
    \item The alignment exceeds simple lexical and template baselines
        (Study B), justifying the compositional-structure interpretation
        rather than the shared-language interpretation.
    \item The alignment is robust to held-out probes and generalizes
        across domains (Study C), justifying generalization claims.
    \item The alignment survives substantial training-distribution
        disjointness (Study D), justifying the substrate-independence
        reading rather than the corpus-residue reading.
    \item Cross-model interventions predicted by the alignment have
        the predicted causal effects (Study E), justifying the
        mechanistic interpretation rather than the correspondence
        interpretation.
    \item Methodological choices have been swept and the result is
        robust (Study F), satisfying ordinary peer-review robustness
        expectations.
    \item Pre-registration is externally verifiable (Study G),
        meeting current open-science standards.
\end{itemize}

We treat that subsequent paper as the work the present pilot
motivates. The present paper is the first step.
